\section{引言：什么是 Harness Engineering}

2026 年 2 月，OpenAI 发表了一篇题为 ``Harness engineering: leveraging Codex in an agent-first world'' 的文章，描述了一个团队如何在\textbf{完全不手动编写代码}的前提下，利用 AI agent 维护一个超过 100 万行代码的大型应用。这一实践持续了 5 个月，产出了一个真实的产品。Thoughtworks 的 Distinguished Engineer Birgitta B\"{o}ckeler 在 Martin Fowler 的博客上撰文，对这篇文章进行了深入的架构视角分析。

\begin{importantbox}{Harness 的核心定义}
``Harness''（直译为``缰绳''或``线束''）在这里指的是一整套\textbf{工具、实践和约束}，用于在 AI agent 驱动的软件开发中保持代码质量和可维护性。它不是单一工具，而是一个\textbf{系统性的控制框架}，结合了确定性规则和 LLM 驱动的检查机制。
\end{importantbox}

值得注意的是，``harness'' 这个术语可能受到了 Mitchell Hashimoto 近期博客文章的启发------OpenAI 的文章标题包含这个词，但正文中仅出现了一次。无论来源如何，Birgitta 认为 ``harness'' 是一个非常恰当的词，用来描述我们可以用来``约束'' AI agent 的工具和实践。

\begin{knowledgebox}{文章背景}
本文是 Birgitta B\"{o}ckeler 在 Martin Fowler 博客上 ``Exploring Gen AI'' 系列的一部分。Birgitta 拥有超过 20 年的软件开发、架构和技术领导经验，目前专注于 AI 辅助交付（AI-assisted delivery）领域。她的分析融合了丰富的工程实践经验与对 AI 工具的深度思考。
\end{knowledgebox}

\subsection{本章小结}

Harness engineering 是一种新兴的工程实践，旨在通过系统性的工具和约束来管控 AI agent 的代码生成行为。它的出现标志着 AI 辅助编程正从``让 AI 生成任意代码''的阶段，走向``有控制地利用 AI 进行大规模软件维护''的阶段。

\newpage

